% Owner: classification_e. Root and boundary_d full-source referee PASS.
\begin{corollary}[E.2.24 (Proved): separation for continuous complex structures]
Let $\mathcal I(S^6)$ be the space of smooth integrable
almost-complex structures on the fixed smooth sphere,
with its $C^\infty$ topology.  For a fixed coprime pair
$(3k,4l)$, let $\mathcal I_{k,l}$ consist of the structures
whose complex manifolds are biholomorphic to constructed
members with that pair.  Then $\mathcal I_{k,l}$ is open
and closed in $\mathcal I(S^6)$.

Consequently a continuous family of such structures
over any connected topological parameter space cannot
contain constructed members with two different pairs.
The same separation holds in the quotient spaces by
$\operatorname{Diff}_0(S^6)$ or by all orientation-preserving
diffeomorphisms, with their quotient topologies.
This statement does not identify the number of marked
connected components belonging to a single pair.
\end{corollary}
\begin{proof}
We use the classical Kuranishi local-completeness
theorem at an arbitrary $J_0\in\mathcal I(S^6)$,
including an initially unconstructed structure.
The complex orientation is locally constant, so the
oriented formulation applies on each such neighbourhood.
The precise $C^\infty$ version used here is
\citet[Section~3.2, Theorem~3.4 and Remark~3.6,
pp.~898--899]{Meersseman2019E}.  It supplies a connected
open neighbourhood $\mathcal U$ of $J_0$, a
finite-dimensional Kuranishi analytic space $K$
containing $J_0$, and a continuous retraction
$\Xi:\mathcal U\to K$ such that every $J\in\mathcal U$
is obtained from $\Xi(J)$ by a diffeomorphism isotopic
to the identity.  In that theorem the local product
map is first constructed in Sobolev completions;
elliptic regularity gives precisely this restriction
to smooth operators.  The topology on these operators
is the $C^\infty$ topology of
\cite[Section~2.1, pp.~883--884]{Meersseman2019E}.

The Kuranishi space carries its proper smooth
holomorphic deformation family.  We can take $K$
connected: in the stated local theorem it is the
continuous retract of the connected $\mathcal U$.
If necessary replace $K$ by its reduction, which
changes neither its points nor the complex manifolds
in its fibres.  Suppose $\mathcal U$ contains some
$J\in\mathcal I_{k,l}$.  Its Kuranishi representative
$\Xi(J)$ has the same complex manifold up to
biholomorphism.  E.2.23 applied to the connected
Kuranishi family therefore makes every point of
$K$ a member of the same pair.  Since all points
of $\mathcal U$ are biholomorphic to their
Kuranishi representatives, it follows that
$\mathcal U\subset\mathcal I_{k,l}$.

Thus at every $J_0$ the neighbourhood $\mathcal U$
is either contained in $\mathcal I_{k,l}$ or
disjoint from it.  This proves both openness and
closedness.  The inverse image of this set under
a continuous map from a connected parameter space
is then empty or the whole space, proving the
continuous-family assertion.  Finally
$\mathcal I_{k,l}$ is saturated under the indicated
diffeomorphism groups, because its definition uses
only the biholomorphism class.  Its image in either
quotient is open and closed by the definition of
the quotient topology.  No assertion that the
saturated set itself is connected is needed.
\end{proof}
